\documentclass[11pt,addpoints,answers]{exam}

\newcommand{\hwnum}{4}
\newcommand{\duedate}{September 30}

\RequirePackage{microtype}
\RequirePackage{mathtools}
\RequirePackage{amsthm}
\RequirePackage{amssymb}
\RequirePackage{xspace}
\RequirePackage[shortlabels]{enumitem}
\RequirePackage{xcolor}
\RequirePackage{hyperref}
\RequirePackage[capitalize,nameinlink,noabbrev]{cleveref} % must load after hyperref
\usepackage{fvextra} % To use Verbatim
\usepackage{diagbox} % For diagbox
\RequirePackage[boxed]{algorithm}
\RequirePackage[noend]{algpseudocode}
\RequirePackage{tikz}
\usetikzlibrary{arrows,automata,positioning}

\hypersetup{breaklinks=true,
    colorlinks=true,
    linkcolor=blue,
    filecolor=blue,
    citecolor=blue,
    urlcolor=blue}

\algrenewcommand{\algorithmiccomment}[1]{\texttt{//} #1}
\algrenewcommand\algorithmicrequire{\textbf{Input:}}
\algrenewcommand\algorithmicensure{\textbf{Output:}}


% allow cleveref to label and reference enumerables defined in the exam class.
% these automatically define corresponding \Crefname as well.

% https://tex.stackexchange.com/questions/126020/cleveref-doesnt-use-correct-capitalized-name-if-used-with-amsthm
\makeatletter
\if@cref@capitalise
\crefname{question}{Question}{Questions}
\Crefname{partno}{Part}{Parts}
\crefname{subpart}{Subpart}{Subparts)}
\crefname{subsubpart}{Subsubpart}{Subsubparts}
\else
\crefname{question}{question}{questions}
\crefname{partno}{part}{parts}
\crefname{subpart}{subpart}{subparts}
\crefname{subsubpart}{subsubpart}{subsubparts}
\fi
\makeatother

% numeric sets in "blackboard" font

\newcommand{\N}{\mathbb{N}}
\newcommand{\Z}{\mathbb{Z}}
\newcommand{\R}{\mathbb{R}}
\newcommand{\Q}{\mathbb{Q}}
\newcommand{\C}{\mathbb{C}}
\newcommand{\Prop}{\mathbb{P}}

% paired delimiters

\DeclarePairedDelimiter\abs{\lvert}{\rvert}
\DeclarePairedDelimiter\length{\lVert}{\rVert}
\DeclarePairedDelimiter\norm{\lVert}{\rVert}
\DeclarePairedDelimiter\parens{(}{)}
\DeclarePairedDelimiter\tuple{(}{)}
\DeclarePairedDelimiter\brackets{[}{]}
\DeclarePairedDelimiter\floor{\lfloor}{\rfloor}
\DeclarePairedDelimiter\ceil{\lceil}{\rceil}
\DeclarePairedDelimiter\round{\lfloor}{\rceil}
\DeclarePairedDelimiter\set{\{}{\}}
\DeclarePairedDelimiter\inner{\langle}{\rangle}

\newcommand{\bit}{\set{0,1}}

% asymptotics

\DeclareMathOperator{\Otil}{\tilde{O}}
\DeclareMathOperator{\poly}{poly}
\DeclareMathOperator{\polylog}{polylog}
\DeclareMathOperator{\negl}{negl}

% algorithms

\newcommand{\algo}[1]{\textsc{#1}}

\newcommand{\memo}{\text{memo}}
\newcommand{\tabl}{\text{table}}
\newcommand{\backtrack}{\text{backtrack}}

\newcommand{\ALG}{\text{ALG}}
\newcommand{\OPT}{\text{OPT}}
\newcommand{\weight}{\text{weight}}
\newcommand{\val}{\text{value}}

% computability

% named language
\newcommand{\lang}[1]{L_{\text{#1}}}
% computational problem
\newcommand{\cproblem}[1]{\ensuremath{\text{#1}}\xspace}
% class of languages
\newcommand{\class}[1]{\ensuremath{\mathsf{#1}}\xspace}

\newcommand{\qst}{q_{\text{start}}}
\newcommand{\qacc}{q_{\text{acc}}}
\newcommand{\qrej}{q_{\text{rej}}}

\newcommand{\Lbarber}{\lang{BARBER}}
\newcommand{\atm}{\lang{ACC}}
\newcommand{\htm}{\lang{HALT}}
\newcommand{\ehtm}{\lang{$\varepsilon$-HALT}}
\newcommand{\eqtm}{\lang{EQ}}
\newcommand{\etm}{\lang{$\emptyset$}}
\newcommand{\epstm}{\lang{$\set{\varepsilon}$}}
\newcommand{\Lprop}{\lang{$\Prop$}}
\newcommand{\LSigmastar}{\lang{$\Sigma^*$}}

% complexity

\newcommand{\yes}{\ensuremath{\text{YES}}}
\newcommand{\no}{\ensuremath{\text{NO}}}

\newcommand{\DTIME}{\class{DTIME}}
\renewcommand{\P}{\class{P}}
\newcommand{\NP}{\class{NP}}
\newcommand{\NPH}{\class{NPH}}
\newcommand{\NPC}{\class{NPC}}
\newcommand{\coNP}{\class{coNP}}

\newcommand{\MAZE}{\cproblem{MAZE}}
\newcommand{\PALINDROME}{\cproblem{PALINDROME}}
\newcommand{\TSP}{\cproblem{TSP}}
\newcommand{\SAT}{\cproblem{SAT}}
\newcommand{\CSAT}{\cproblem{CSAT}}
\newcommand{\TSAT}{\cproblem{3SAT}}
\newcommand{\VC}{\cproblem{VERTEX-COVER}}
\newcommand{\DS}{\cproblem{DOMINATING-SET}}
\newcommand{\SC}{\cproblem{SET-COVER}}
\newcommand{\HC}{\cproblem{HAMCYCLE}}
\newcommand{\HP}{\cproblem{HAMPATH}}
\newcommand{\AHP}{\cproblem{ANYHAMPATH}}
\newcommand{\IS}{\cproblem{IS}}
\newcommand{\CLIQUE}{\cproblem{CLIQUE}}
\newcommand{\SSUM}{\cproblem{SUBSET-SUM}}
\newcommand{\KNAPSACK}{\cproblem{KNAPSACK}}
\newcommand{\MAXCUT}{\cproblem{MAX-CUT}}

% randomness

\DeclareMathOperator*{\Var}{Var}
\DeclareMathOperator*{\Ex}{\mathbb{E}}

\newcommand{\RP}{\class{RP}}
\newcommand{\coRP}{\class{coRP}}
\newcommand{\BPP}{\class{BPP}}
\newcommand{\ZPP}{\class{ZPP}}
\newcommand{\BQP}{\class{BQP}}

%%% misc

\newcommand{\eps}{\varepsilon}

%%% theorems

\theoremstyle{plain}            % following are "theorem" style

\newtheorem{theorem}{Theorem}
\newtheorem{lemma}[theorem]{Lemma}
\newtheorem{corollary}[theorem]{Corollary}
\newtheorem{proposition}[theorem]{Proposition}
\newtheorem{claim}[theorem]{Claim}
\newtheorem{fact}[theorem]{Fact}
\newtheorem{openproblem}[theorem]{Open Problem}

\theoremstyle{definition}       % following are def style

\newtheorem{definition}[theorem]{Definition}
\newtheorem{conjecture}[theorem]{Conjecture}
\newtheorem{protocol}[theorem]{Protocol}
\newtheorem{exercise}[theorem]{Exercise}

\theoremstyle{remark}           % following are remark style

\newtheorem{example}[theorem]{Example}
\newtheorem{remark}[theorem]{Remark}
\newtheorem{note}[theorem]{Note}

%%% for homework and section notes

\newcommand{\commonheader}[2]{
    \pagestyle{headandfoot}
    \setlength{\headheight}{26pt}
    \setlength{\headsep}{16pt}

    \header
        {\small{\textbf{EECS 376: Foundations of Computer Science}} \\ \footnotesize{\textbf{University of Michigan, Fall 2026}}}
        {#1}
        {#2}

    \firstpageheadrule
    \runningheadrule

    \footer
        {}
        {\thepage}
        {}
}

\newcommand{\hwheader}{
    \commonheader
        {\Large \textbf{Homework \hwnum}}
        {\small \textbf{Due 8:00pm, \duedate\\ {\tiny(accepted until 10:00 pm, no credit after)}}}
}

\newcommand{\hwslnheader}{
    \commonheader
    	{}
        {\Large \textbf{Solutions to Homework \hwnum}}
    \printanswers
}

\newcommand{\notesheader}{
    \commonheader
    	{}
        {\Large \textbf{Discussion Notes \sectionnum}}
}

\newcommand{\practiceheader}{
    \commonheader
    	{}
        {\Large \textbf{Discussion Worksheet \sectionnum}}
}

\newcommand{\practiceslnheader}{
    \commonheader
    	{}
        {\Large \textbf{Solutions to Discussion Worksheet \sectionnum}}
}

\newcommand{\reviewheader}{
    \commonheader 
    \smallskip
    	{}
        {\Large \textbf{Midterm Review Notes}}
}

\newcommand{\hwpreface}{

\noindent This homework has \numquestions\ questions, for a total of \numpoints\ points and \numbonuspoints\ extra-credit points.

\noindent Unless otherwise stated, each question requires \emph{clear}, \emph{logically correct}, and \emph{sufficient} justification to convince the reader.

\noindent We strongly encourage you to typeset your solutions in \LaTeX.

\noindent If you collaborated with someone, you must state their name(s). You must \emph{write your own solution} for all problems and \emph{may not use any other student’s write-up or any AI tools}.
}

\newcommand{\hint}[1]{
\emph{Hint}: #1
}

% exam class setup
\pointsinmargin
\pointpoints{pt}{pts}
\bonuspointpoints{EC pt}{EC pts}
\marginpointname{ \points}
\marginbonuspointname{ \bonuspoints}

\usepackage{tikz}
\usetikzlibrary{arrows.meta,positioning}

\hwslnheader

\begin{document}

\hwpreface

\begin{questions}

  \addtocounter{question}{-1}
  \question \textbf{Before you start; before you submit.}

  If applicable, state the name(s) and uniqname(s) of your collaborator(s).

  \begin{solution}
  No collaborators.
  \end{solution}

\question[10] \textbf{Self assessment.} \nopagebreak
  Carefully read and understand the posted solutions to the previous homework.
  Identify one part for which your own solution has the most room for improvement (e.g., has unsound reasoning, doesn’t show what was required, could be significantly clearer or better organized, etc.).
  Copy or screenshot this solution, then in a few sentences, explain what was deficient and how it could be fixed.

  (Alternatively, if you think one of your solutions is significantly \emph{better} than the posted one, copy it here and explain why you think it is better.)

  If you didn't do last week's homework, choose a problem from it that looks challenging to you, and in a few sentences, explain the key ideas behind its solution in your own words.

  \begin{solution}
  I chose the minimum-subarray-product problem from the previous homework as a
  useful example to reflect on.  The key observation is that every nonempty
  subarray ending at position $i$ is either the one-element subarray $[A_i]$,
  or is obtained by appending $A_i$ to a subarray ending at $i-1$.  Thus, if
  $E_i$ denotes the smallest product of a nonempty subarray ending at $i$,
  then
  \[
      E_i=\min\{A_i,\,E_{i-1}A_i\}.
  \]
  The answer is $\min_i E_i$.  A solution can be improved by explicitly
  separating the local state $E_i$ from the global answer, stating the base
  case $E_1=A_1$, and explaining why the two cases above exhaust all
  subarrays ending at $i$.  That organization makes both correctness and the
  $O(n)$ running time immediate.
  \end{solution}

\question \textbf{Fixed-color-sequence paths.}

Let $G=(V,E)$ be an undirected graph with $n$ vertices and $m$ edges, where
$m\ge n$.  Each edge is colored black or white.  Given vertices $s,t$ and a
desired color sequence $(c_1,\ldots,c_k)$, we must determine whether there is
a length-$k$ path from $s$ to $t$ with exactly this sequence.

  \begin{parts}
    \part[10]
    \begin{solution}
    Define the DP state $F(i,v)$ to be true exactly when there is a length-$i$
    path from $s$ to $v$ whose edge colors are $(c_1,\ldots,c_i)$.  The base
    cases are
    \[
      F(0,s)=\text{true},\qquad F(0,v)=\text{false}\quad(v\ne s).
    \]
    For $i\ge1$,
    \[
      F(i,v)=\bigvee_{\substack{u\in\operatorname{Adj}(v)\\
                    \operatorname{color}(u,v)=c_i}}F(i-1,u).
    \]
    The desired answer is $F(k,t)$.

    This recurrence has the needed optimal-substructure property: the prefix
    before the last edge must itself be a valid solution to the smaller
    length-$i-1$ subproblem.  To justify the recurrence, consider the last edge
    of a length-$i$ path
    ending at $v$.  It must come from some neighbor $u$ whose edge has color
    $c_i$, and the preceding edges must form a valid length-$(i-1)$ path from
    $s$ to $u$.  This proves that every valid path is represented by one true
    term in the disjunction.  Conversely, any true term extends a valid
    length-$(i-1)$ path by the qualifying edge $(u,v)$, producing a valid
    length-$i$ path.  The base cases describe the unique length-zero path.
    \end{solution}

    \part[6]
    \begin{solution}
    \begin{algorithm}[H]
      \caption{FixedColorPath$(G,s,t,c_1,\ldots,c_k)$}
      \begin{algorithmic}[1]
        \ForAll{$v\in V$}
          \State{$F[0,v]\gets\text{false}$}
        \EndFor
        \State{$F[0,s]\gets\text{true}$}
        \For{$i\gets1$ \textbf{to} $k$}
          \ForAll{$v\in V$}
            \State{$F[i,v]\gets\text{false}$}
            \ForAll{$u\in\operatorname{Adj}(v)$}
              \If{$\operatorname{color}(u,v)=c_i$ \textbf{and} $F[i-1,u]$}
                \State{$F[i,v]\gets\text{true}$}
              \EndIf
            \EndFor
          \EndFor
        \EndFor
        \State{\Return{$F[k,t]$}}
      \end{algorithmic}
    \end{algorithm}

    This pseudocode evaluates the recurrence by bottom-up table filling.  For
    a fixed $i$, the inner adjacency-list scans examine a total of
    $2m$ entries, because the graph is undirected and each edge occurs in two
    adjacency lists.  The work for one layer is therefore $O(m+n)=O(m)$,
    since $m\ge n$.  There are $k$ layers, so the total running time is
    $O(km)$.  The table uses $O(kn)$ space; keeping only the previous layer
    would reduce the space to $O(n)$ without changing the time.
    \end{solution}
  \end{parts}

\question \textbf{Back to college}
\nopagebreak

  Let $x_0=0$ and $x_i=\sum_{r=1}^i d_r$ be the total distance from town $0$
  to town $i$.

  \begin{parts}
    \part[7]
    \begin{solution}
    Starting immediately after a refueling at town $p$, the greedy choice is
    the farthest town $q$ that can be reached on the current full pack:
    \[
      q=\max\{r\in\{p+1,\ldots,n\}:x_r-x_p\le R\}.
    \]
    If $q=n$, no further refueling is needed.  Otherwise refuel at $q$ and
    repeat.

    \begin{algorithm}[H]
      \caption{GreedyRefuel$(R,D)$}
      \begin{algorithmic}[1]
        \State{$x[0]\gets0$}
        \For{$i\gets1$ \textbf{to} $n$}
          \State{$x[i]\gets x[i-1]+d_i$}
        \EndFor
        \State{$S\gets\emptyset$; $p\gets0$}
        \While{$p<n$}
          \State{$q\gets p$}
          \While{$q<n$ \textbf{and} $x[q+1]-x[p]\le R$}
            \State{$q\gets q+1$}
          \EndWhile
          \If{$q=n$}
            \State{\textbf{break}}
          \EndIf
          \State{$S\gets S\cup\{q\}$}
          \State{$p\gets q$}
        \EndWhile
        \State{\Return{$S$}}
      \end{algorithmic}
    \end{algorithm}

    The prefix sums take $O(n)$ time.  The pointer $q$ only moves forward,
    so over all iterations the two pointers together also take $O(n)$ time.
    The algorithm uses $O(n)$ space for the prefix sums, or $O(1)$ additional
    space if the distances are scanned with a suitable running total.
    \end{solution}

    \part[10]
    \begin{solution}
    Use the notation from the prompt, and let $p=i_{\mathrm{diff}}-1$ be the
    previous greedy refueling town.  Because $\OPT$ is valid, it must have a
    next refueling town after $p$ unless it can reach town $n$ directly from
    $p$.  The latter case is impossible: the greedy algorithm would then have
    chosen $n$ instead of $i_{\mathrm{diff}}$.  Let $r$ be the first
    refueling town of $\OPT$ after $p$.  Since $i_{\mathrm{diff}}$ is the
    farthest reachable town from $p$, $r\le i_{\mathrm{diff}}$; and since
    $i_{\mathrm{diff}}\notin\OPT$, we have $r<i_{\mathrm{diff}}$.
    Let $j$ be the last town in $\OPT$ that occurs after $p$ and before
    $i_{\mathrm{diff}}$.  Thus $j$ exists and $j<i_{\mathrm{diff}}$.

    Let $q$ be the next refueling town of $\OPT$ after $j$, or let $q=n$ if
    there is no such refueling town.  Because $\OPT$ is valid,
    \[
       x_q-x_j\le R.
    \]
    The greedy choice from $p$ satisfies
    $x_{i_{\mathrm{diff}}}-x_p\le R$.  Therefore, replacing $j$ by
    $i_{\mathrm{diff}}$ preserves feasibility on both affected segments:
    \[
       x_{i_{\mathrm{diff}}}-x_p\le R,
       \qquad
       x_q-x_{i_{\mathrm{diff}}}
       \le x_q-x_j\le R,
    \]
    where the second inequality uses $j<i_{\mathrm{diff}}\le q$.
    All other consecutive refueling segments are unchanged.  Thus
    \[
       \OPT'=(\OPT\setminus\{j\})\cup\{i_{\mathrm{diff}}\}
    \]
    is a valid solution.

    The exchange removes exactly one town and adds exactly one town, so
    $|\OPT'|=|\OPT|$.  Since $\OPT$ is optimal, no valid solution can use
    fewer refueling towns, so $\OPT'$ is also optimal.  It contains the
    greedy prefix through $i_{\mathrm{diff}}$, as required.
    \end{solution}

    \part[4]
    \begin{solution}
    The exchange argument shows that each greedy choice is safe: whenever an
    optimal solution first
    disagrees with the greedy solution, it can be modified into another
    optimal solution that agrees with one more greedy choice.  Repeating this
    argument makes an optimal solution contain the entire greedy sequence.
    Since the greedy sequence itself is valid, and an optimal solution cannot
    have fewer refueling towns than a valid greedy solution containing its
    whole sequence, the two sequences have the same size.  Hence the greedy
    algorithm is optimal.
    \end{solution}

    \part[5]
    \begin{solution}
    The claim is false.  Consider
    \[
       R=10,\qquad D=[6,4,6],\qquad C=[3,8,100].
    \]
    The town positions are $x_0=0,x_1=6,x_2=10,x_3=16$.  The greedy
    algorithm chooses town $2$, because it is the farthest town reachable from
    town $0$ on the initial pack.  It therefore pays cost $c_2=8$.

    A better solution refuels at town $1$: after traveling $6$ units, the new
    pack covers the remaining distance $4+6=10$, so the car reaches town $3$
    with total refueling cost $c_1=3$.  Thus the greedy solution is not
    minimum-cost.

    The exchange proof from the previous part fails because exchanging town
    $j$ for the later town $i_{\mathrm{diff}}$ preserves the number of
    refuelings, but it need not preserve or decrease their monetary cost.
    \end{solution}
  \end{parts}

\question \textbf{Dynamic programming for dynamic shortest paths.} \nopagebreak

Let $v_{\mathrm{new}}$ be the added vertex.

  \begin{parts}
    \part[10]
    \begin{solution}
    For every $u\in V$,
    \[
      D_{\mathrm{new}}(v_{\mathrm{new}},u)
       =\min_{(v_{\mathrm{new}},y)\in E_{\mathrm{new}}}
          \bigl(w(v_{\mathrm{new}},y)+D(y,u)\bigr),
      \tag{1}
    \]
    and
    \[
      D_{\mathrm{new}}(u,v_{\mathrm{new}})
       =\min_{(y,v_{\mathrm{new}})\in E_{\mathrm{new}}}
          \bigl(D(u,y)+w(y,v_{\mathrm{new}})\bigr).
      \tag{2}
    \]
    Empty minima are interpreted as $+\infty$.

    To prove (1), take a shortest path from $v_{\mathrm{new}}$ to $u$.
    Because there is no negative-weight cycle, a shortest path can be chosen
    simple, so after its first edge it never returns to $v_{\mathrm{new}}$.
    If that first edge is $(v_{\mathrm{new}},y)$, the remaining path lies
    entirely in the old graph and has length at least $D(y,u)$.  This proves
    that every such path has length at least the right-hand side.  Conversely,
    choosing any outgoing edge and concatenating it with an old shortest path
    from $y$ to $u$ gives a valid path of the displayed length, proving
    equality.  The proof of (2) is the same argument applied to the final
    edge entering $v_{\mathrm{new}}$.

    If $d^+$ is the number of outgoing new edges, evaluating (1) for all
    $u$ takes $O(nd^+)$ time.  If $d^-$ is the number of incoming new edges,
    evaluating (2) takes $O(nd^-)$ time.  Since $d^++d^-=O(n)$ in a simple
    graph, the total is $O(n^2)$.
    \end{solution}

    \part[8]
    \begin{solution}
    For every $u,v\in V$,
    \[
      D_{\mathrm{new}}(u,v)=\min\left\{
        D(u,v),\;
        D_{\mathrm{new}}(u,v_{\mathrm{new}})
        +D_{\mathrm{new}}(v_{\mathrm{new}},v)
      \right\}.
      \tag{3}
    \]

    This is the same shortest-path recurrence idea used in dynamic
    programming for graphs.  A shortest path from $u$ to $v$ either avoids
    $v_{\mathrm{new}}$, in
    which case its length is at least $D(u,v)$, or uses it.  In the latter
    case, split the path at its occurrence of $v_{\mathrm{new}}$.  The first
    part has length at least
    $D_{\mathrm{new}}(u,v_{\mathrm{new}})$ and the second part has length at
    least $D_{\mathrm{new}}(v_{\mathrm{new}},v)$, so its length is at least
    the second term of (3).  This proves that every path has length at least
    the right-hand side.  Conversely, each term in (3) is realized by a
    valid old-graph path or by concatenating two valid paths through the new
    vertex, so the minimum is attainable.  If a concatenation repeats a
    vertex, it is still a valid walk; since there is no negative cycle, any
    repeated nonnegative cycle can be removed without increasing its length.

    The two quantities in (3) have already been computed in the previous
    part, so each pair $(u,v)$ takes $O(1)$ time.  There are $n^2$ pairs,
    giving $O(n^2)$ additional time and therefore $O(n^2)$ total time.
    \end{solution}
  \end{parts}

\question \textbf{MST mashup.} \nopagebreak

  \begin{parts}
    \part[5]
    \begin{solution}
    The problematic implication is the final one: \textquotedblleft Kruskal has one output; therefore that output is the only minimum spanning tree.\textquotedblright\ Correctness
    of Kruskal says that its output is an MST, not that every MST must be the
    output of a particular deterministic algorithm.  An algorithm can make a
    single deterministic choice even when several different optimal solutions
    exist.  Thus uniqueness of Kruskal's execution does not imply uniqueness
    of the set of all MSTs.
    \end{solution}

    \part[8]
    \begin{solution}
    Suppose for contradiction that two distinct spanning trees $T_1$ and
    $T_2$ are both minimum spanning trees.  Choose an edge $e$ of minimum
    weight in the symmetric difference $T_1\triangle T_2$, and without loss
    of generality assume $e\in T_1\setminus T_2$.

    Remove $e$ from $T_1$.  This creates two connected components, say $A$
    and $B$.  Consider the unique $T_2$-path between the endpoints of $e$.
    Its endpoints lie on opposite sides of $(A,B)$, so the path contains an
    edge $f$ crossing this cut.  The edge $f$ cannot belong to $T_1-e$, since
    that graph has no edge crossing the cut, and it is not $e$ because
    $e\notin T_2$.  Hence $f\in T_2\setminus T_1$, so $f$ also belongs to
    the symmetric difference.

    All edge weights are distinct, and $e$ was chosen as the minimum-weight
    edge in the symmetric difference.  Therefore $w(e)<w(f)$.  Now
    Adding $e$ to $T_2$ creates the unique cycle formed by $e$ and the
    $T_2$-path just considered; $f$ lies on this cycle.  Therefore
    $T_2-f+e$ is a spanning tree.  Its weight is
    \[
       w(T_2)-w(f)+w(e)<w(T_2),
    \]
    contradicting the minimality of $T_2$.  Thus a graph with distinct edge
    weights has at most one MST; since a spanning tree exists in the setting
    of the claim, its MST is unique.
    \end{solution}
  \end{parts}

\question \textbf{Greedy knapsack filling.} \nopagebreak

  \begin{parts}
    \part[2]
    \begin{solution}
    Every feasible 0--1 solution is also a feasible fractional solution by
    choosing each fraction $p_i$ to be either $0$ or $1$.  The fractional
    problem maximizes over a superset of the 0--1 feasible solutions, so its
    optimal value is at least as large.
    \end{solution}

    \part[5]
    \begin{solution}
    The proposed algorithm is not correct.  Consider capacity $C=10$ and
    three items:
    \[
      (w_1,v_1)=(100,100),\qquad
      (w_2,v_2)=(1,1),\qquad
      (w_3,v_3)=(10,90).
    \]
    \algo{ValueGreedy} first chooses item $1$, since it has value $100$.
    It takes a $1/10$ fraction and obtains value $10$.

    \algo{WeightGreedy} first takes all of item $2$, using one unit of
    capacity.  It then takes a $9/10$ fraction of item $3$, obtaining total
    value $1+81=82$.  Thus the better of the two greedy outputs has value
    $82$.

    However, taking all of item $3$ uses exactly the capacity and gives value
    $90$, which is optimal: item $3$ alone has value $90$, and no feasible
    solution can use more than the maximum value-per-weight item $3$ on all
    ten units of capacity.  Therefore the proposed algorithm can return a
    suboptimal answer.
    \end{solution}

    \part[10]
    \begin{solution}
    Let $r_i=v_i/w_i$ and assume $r_1\ge r_2\ge\cdots\ge r_n$.  Let
    $p_i$ be the fractions selected by \algo{RelativelyGreedy}, and let
    $q_i$ be the fractions of an arbitrary optimal solution $\OPT$.
    Choose the smallest index $j$ for which $q_j\ne p_j$.  The fractions for
    indices below $j$ already agree, so let
    \[
       C'=C-\sum_{i<j}p_iw_i
    \]
    be the capacity remaining before item $j$.

    First suppose $p_j<1$.  Then the greedy algorithm fills the remaining
    capacity with item $j$, so $p_jw_j=C'$ (the case $p_j=0$ cannot be a
    first mismatch with a feasible solution unless there is no remaining
    capacity).  Feasibility of $\OPT$ implies
    \[
       q_jw_j+\sum_{i>j}q_iw_i\le C'=p_jw_j,
    \]
    and hence $q_j\le p_j$.  Let $\delta=p_j-q_j>0$.  There is at least
    $\delta w_j$ of capacity available either as unused capacity or as weight
    currently assigned to later items.  Add a $\delta$-fraction of item $j$,
    and remove a total weight $\delta w_j$ from later selected items, using
    unused capacity first and then reducing later fractions as needed.  This
    yields a feasible solution $\OPT'$ with the required fraction of item $j$.

    Now suppose $p_j=1$.  Then $C'\ge w_j$.  Since $q_j\ne1$ and
    $0\le q_j\le1$, let $\delta=1-q_j>0$.  The quantity of unused capacity
    plus the weight assigned to later items is
    \[
       C'-q_jw_j\ge w_j-q_jw_j=\delta w_j,
    \]
    so the same exchange is possible: add a $\delta$-fraction of item $j$
    and remove total weight $\delta w_j$ from unused capacity and later
    items.

    In either case, every removed fraction belongs to an item $i>j$, whose
    relative value satisfies $r_i\le r_j$.  The value added is
    $\delta w_jr_j$.  The value removed from later items is at most
    $\delta w_jr_j$, and unused capacity removes no value.  Consequently
    \[
       \val(\OPT')\ge\val(\OPT).
    \]
    Since $\OPT$ was optimal, $\OPT'$ is also optimal.  It agrees with the
    greedy fractions through index $j$.  Repeating this exchange for the next
    disagreement eventually transforms an optimal solution into the greedy
    solution.  Therefore \algo{RelativelyGreedy} is correct.
    \end{solution}
  \end{parts}

\bonusquestion \textbf{Optional extra credit: Greedily independent.}
\nopagebreak

  \begin{parts}
    \bonuspart[1]
    \begin{solution}
    Use vertices $\{1,2,3,4,5,6\}$ and edges
    \[
      \{(1,2),(1,3),(1,4),
        (2,5),(2,6),(3,5),(3,6),(4,5),(4,6),(5,6)\}.
    \]
    Vertex $1$ has degree $3$, vertices $2,3,4$ have degree $3$, and
    vertices $5,6$ have degree $5$.  Therefore the tie-breaking rule selects
    vertex $1$ first.  It removes vertices $2,3,4$, after which the algorithm
    selects one of $5,6$, returning an independent set of size $2$.  But
    $\{2,3,4\}$ is an independent set of size $3$, so the output is not
    maximum.
    \end{solution}

    \bonuspart[2]
    \begin{solution}
    Consider the graph remaining immediately before the algorithm selects
    $v_{\mathrm{diff}}$.  It is a forest, because deleting vertices and edges
    from a tree leaves a forest.  Every nonempty forest has a vertex of
    degree at most $1$, so the minimum-degree choice satisfies
    $\deg(v_{\mathrm{diff}})\le1$ in the remaining forest.

    All previously chosen greedy vertices belong to $\OPT$ by the setup, and
    the remaining vertex $v_{\mathrm{diff}}$ is not adjacent to any of them;
    otherwise it would already have been deleted.  If
    $\deg(v_{\mathrm{diff}})=0$, then $v_{\mathrm{diff}}$ has no neighbor in
    $\OPT$, so we could add it to $\OPT$, contradicting maximality of
    $\OPT$.  Thus it has exactly one remaining neighbor, call it $u$, and
    $u\in\OPT$.

    Define
    \[
       \OPT'=(\OPT\setminus\{u\})\cup\{v_{\mathrm{diff}}\}.
    \]
    The only possible neighbor of $v_{\mathrm{diff}}$ in the remaining
    forest is $u$.  Any vertex deleted earlier is either a previously chosen
    greedy vertex (which is in $\OPT$ and is not adjacent to
    $v_{\mathrm{diff}}$) or a neighbor of one of those vertices (and hence
    cannot belong to the independent set $\OPT$).  Therefore adding
    $v_{\mathrm{diff}}$ after removing $u$ creates no adjacency inside
    $\OPT'$.  The set is independent and has the same cardinality as $\OPT$,
    so it is also a maximum independent set.  It contains the greedy prefix
    through $v_{\mathrm{diff}}$.

    Repeating the exchange for each first disagreement produces a maximum
    independent set containing the entire greedy output.  Since the greedy
    output itself is independent, it must have the same size as a maximum
    independent set and is therefore maximum.
    \end{solution}
  \end{parts}

\end{questions}
\end{document}
